\documentclass[../main.tex]{subfiles}
\graphicspath{{\subfix{../Images/}}}

\begin{document}
\onecolumn
% =================================================================
% Section IV -- Problem Formulation (rewritten)
%
% Companion to method_implementation_skeleton.tex.
%
% CONTRACT
%   This section says WHAT the objects are: sets, constraints, the coupling
%   condition, and the quantities used to score a run.
%   It never says HOW anything is computed. Every "how" is a forward
%   reference to Section V (Method) or Section VI (Implementation).
%
% Every equation carries a \label consumed by the Method skeleton.
% Labels exported by this section:
%   eq:onestepcost  eq:actions      eq:transition    eq:vertexconflict
%   eq:swapconflict eq:feasiblestorage  eq:storageupdate  eq:storageobjective
%   eq:task         eq:coupling     eq:servicetime   eq:partition
%   eq:localsubgraph eq:observation eq:potential     eq:commgraph
%   eq:posg         eq:congestion   eq:objective     eq:throughput
%   eq:movementcost eq:entropy      eq:concentration eq:functional
%   eq:decisionproblem  eq:rqformal
% Section labels: sec:problem, sec:pf:env, sec:pf:agents, sec:pf:storage,
%   sec:pf:tasks, sec:pf:controllers, sec:pf:observations, sec:pf:game,
%   sec:pf:coupling, sec:pf:measures, sec:pf:scope
%
% Requires: amsmath, amssymb, graphicx, booktabs
% =================================================================

\section{Problem Formulation}
\label{sec:problem}

This section defines the coupled storage--routing system studied in this thesis.
It proceeds in the order in which a simulation run unfolds: the environment, the
agents that move in it, the storage state, the work that arrives, the controllers
that decide, what a decentralised agent can see, the resulting learning problem,
the loop that closes between the two layers, and the quantities used to score a
run. Each object is described in plain terms before it is written formally.

Throughout, this section states \emph{what} the objects are, not \emph{how} they
are computed. Realisations --- heuristics, algorithms, network architectures and
hyperparameters --- are deferred to Section~\ref{sec:method}, and software
concerns to Section~\ref{sec:implementation}.

% -----------------------------------------------------------------
\subsection{The environment}
\label{sec:pf:env}

The warehouse is a finite graph $G=(V,E)$. Vertices are the discrete positions a
robot can occupy: aisle cells, cross-aisle junctions, storage faces and drop-off
points. Edges are the movements that are physically possible between two adjacent
positions. A rectangular grid is one special case, but the formulation
deliberately does not assume a lattice, because real warehouses contain one-way
aisles, diagonal shortcuts and irregular blocks.

Four roles partition the vertices by what can happen at them.
$V_{\mathrm{mov}} \subseteq V$ are the vertices an agent may occupy.
$V_{\mathrm{str}} \subseteq V$ are storage faces, where items are held and from
which pickups are made. $V_{\mathrm{del}} \subseteq V$ are delivery or drop-off
points, where completed tasks terminate. $V_{\mathrm{ep}} \subseteq
V_{\mathrm{mov}}$ are endpoints: vertices at which an idle agent may remain
without blocking transit. Each edge $e \in E$ carries a traversal cost
$c(e) > 0$, and $d_G(v,w)$ denotes the shortest-path cost from $v$ to $w$.

Figure~\ref{fig:instance} shows a small instance with all four roles marked.

\begin{figure}[t]
  \centering
  \includegraphics[width=\linewidth]{Images/instance.png}
  \caption{A small warehouse instance. Vertex shape encodes role: storage faces,
  delivery points, endpoints and plain transit vertices. A single task
  $\tau_j = (r_j, s_j, g_j)$ is shown as a route from its pickup vertex $s_j$ to
  its delivery vertex $g_j$. The layout is not a lattice: aisle structure, not
  grid geometry, determines which movements exist.}
  \label{fig:instance}
\end{figure}

At each timestep an agent at vertex $v$ either waits or traverses one outgoing
edge. Writing $N^{+}(v) = \{w \in V_{\mathrm{mov}} : (v,w) \in E\}$ for the
movable successors of $v$, the action set at $v$ is
\begin{equation}
  \mathcal{U}(v) = \{\mathrm{wait}\} \cup \{\mathrm{move}(w) : w \in N^{+}(v)\}.
  \label{eq:actions}
\end{equation}
Waiting costs $c_{\mathrm{wait}} \geq 0$ and moving from $v$ to $w$ costs
$c(v,w)$, so the realised one-step cost is
\begin{equation}
  \hat{c}(v,w) =
  \begin{cases}
    c_{\mathrm{wait}}, & v = w, \\
    c(v,w),            & (v,w) \in E.
  \end{cases}
  \label{eq:onestepcost}
\end{equation}
The two cases are needed because a self-loop $(v,v)$ is not an element of $E$, so
$c$ is undefined on a wait and any route cost containing a wait would otherwise be
ill-typed.\footnote{The value of $c_{\mathrm{wait}}$ relative to edge costs
determines whether an agent prefers to yield or to detour; it is therefore a
parameter of the study rather than a constant of the model, and is reported in
Section~\ref{sec:method:training}.}

Some baseline algorithms are complete only on \emph{well-formed} instances. Where
such a baseline is used, the following is imposed as an assumption on the
generated instance rather than derived: at least $m$ non-task endpoints exist,
and between any pair of endpoints there is a path that does not pass through a
third endpoint. Instances are checked against this condition in
Section~\ref{sec:impl:instances}, and the condition is reported only for the
baselines to which it applies.

% -----------------------------------------------------------------
\subsection{Agents and collisions}
\label{sec:pf:agents}

A fleet $\mathcal{A} = \{a_1,\dots,a_m\}$ of homogeneous agents moves on $G$.
Agent $a_i$ occupies $\ell_i(t) \in V_{\mathrm{mov}}$ at time $t$ and selects
$u_i(t) \in \mathcal{U}(\ell_i(t))$. Its realised route over the horizon is
$p_i = (\ell_i(0), \dots, \ell_i(T))$. A single agent's transition is legal when
\begin{equation}
  \ell_i(t+1) = \ell_i(t)
  \quad \vee \quad
  \bigl(\ell_i(t), \ell_i(t+1)\bigr) \in E .
  \label{eq:transition}
\end{equation}

Legality of each agent's own move is not sufficient, because agents share the
infrastructure. Two agents collide if they end a timestep at the same vertex, or
if they exchange positions across a single edge. A joint transition is therefore
\emph{collision-free} when both
\begin{align}
  \ell_i(t+1) &\neq \ell_j(t+1) && \forall i \neq j,
  \label{eq:vertexconflict} \\
  \neg\bigl( \ell_i(t) = \ell_j(t+1) \;\wedge\; \ell_j(t) = \ell_i(t+1) \bigr)
      & && \forall i \neq j .
  \label{eq:swapconflict}
\end{align}
The two conditions are logically independent. Constraint
\eqref{eq:vertexconflict} inspects the successor state alone;
\eqref{eq:swapconflict} inspects the pair (state, successor) and is invisible to
any check performed only at $t+1$. Omitting it permits two agents to pass through
one another in a single-lane aisle. Figure~\ref{fig:conflicts} illustrates both,
together with the feasible alternative in which one agent waits.

\begin{figure}[t]
  \centering
  \includegraphics[width=\linewidth]{Images/conflicts.png}
  \caption{The two forbidden joint transitions and one feasible alternative.
  Panel~(c) shows the resolution the model always admits: waiting is a legal
  action at every vertex, so a collision-free joint action always exists. Which
  agent yields is a decision, not a constraint, and is fixed in
  Section~\ref{sec:method:resolution}.}
  \label{fig:conflicts}
\end{figure}

Constraints \eqref{eq:vertexconflict}--\eqref{eq:swapconflict} are hard: a run in
which they are violated is not a solution to this problem. They are stated here
as feasibility conditions on the realised trace. A learned policy induces a
distribution over joint actions whose support may contain infeasible ones, so
feasibility cannot be delegated to a penalty term in a reward; the mechanism that
enforces it is specified in Sections~\ref{sec:method:rl}
and~\ref{sec:method:resolution} and verified in Section~\ref{sec:impl:validation}.

% -----------------------------------------------------------------
\subsection{Where the items are}
\label{sec:pf:storage}

Items are held at storage faces, and where they are held determines where pickups
happen. Let $\mathcal{K}$ be the set of SKU types. The storage state must be able
to represent split inventory, multiple units, and finite shelf space, so it is
defined as a counting function
\[
  x_t : \mathcal{K} \times V_{\mathrm{str}} \to \mathbb{N}_0 ,
\]
where $x_t(k,v)$ is the number of units of SKU $k$ held at storage vertex $v$ at
time $t$. Let $b_k > 0$ be the capacity consumed by one unit of SKU $k$ and
$\mathrm{cap}(v) > 0$ the capacity of vertex $v$. The set of feasible
configurations is
\begin{equation}
  \mathcal{X} =
  \Bigl\{ x \in \mathbb{N}_0^{|\mathcal{K}| \times |V_{\mathrm{str}}|} :
    \sum_{k \in \mathcal{K}} b_k\, x(k,v) \le \mathrm{cap}(v)
    \quad \forall v \in V_{\mathrm{str}} \Bigr\}.
  \label{eq:feasiblestorage}
\end{equation}
A simpler mapping $\sigma_t : \mathcal{K} \to V_{\mathrm{str}}$ would force one
location per SKU, unit quantity and unlimited shelf space; it is admissible only
if those three restrictions are stated explicitly. Note that
\eqref{eq:feasiblestorage} is a system of $|V_{\mathrm{str}}|$ knapsack
constraints, so capacity couples the placement decisions of different SKUs at a
shared vertex.

The storage layer acts on a slow clock. At epochs
$t \in \{\Delta, 2\Delta, \dots\}$ the configuration is updated by a rule $F$,
\begin{equation}
  x_t = F\bigl(x_{t-\Delta},\, \hat{\rho}_t,\, \hat{\mu}_t,\, \hat{w}_t\bigr)
        \in \mathcal{X} ,
  \label{eq:storageupdate}
\end{equation}
where $\hat{\rho}_t$ estimates per-SKU demand, $\hat{\mu}_t(e)$ estimates
traversal counts per edge, and $\hat{w}_t(e)$ estimates waiting attributed to an
edge or its downstream contention point. All three are computed only from data
available strictly before the epoch, so \eqref{eq:storageupdate} is causal. The
fixed-storage baseline is the identity rule $F_{\mathrm{fix}}(x,\cdot) = x$,
defined directly rather than as the informal limit $\Delta = \infty$.

Because capacity couples SKUs, the update cannot be a per-item greedy score
applied independently. It is a joint constrained decision over $\mathcal{X}$,
\begin{equation}
  x_t \in \arg\min_{x \in \mathcal{X}}
    \bigl[\, \alpha D(x; \hat{\rho}_t)
       + \beta K(x; \hat{\mu}_t, \hat{w}_t)
       + \chi R(x, x_{t-\Delta}) \,\bigr],
  \label{eq:storageobjective}
\end{equation}
where $D$ measures demand-weighted access distance, $K$ measures predicted
congestion exposure, and $R$ penalises reassignment. In words: $D$ says popular
items should be close, $K$ says they should not all sit on the same corridor, and
$R$ says the warehouse should not be rearranged every epoch. The three functional
forms, the weights $\alpha, \beta, \chi$, and the heuristic that actually
searches $\mathcal{X}$ are specified in Section~\ref{sec:method:storage};
\eqref{eq:storageobjective} is an integer program and is not solved to
optimality. Setting $\chi = 0$ removes relocation cost from the model entirely
and is recorded in Section~\ref{sec:pf:scope} as a limitation rather than a
simplification in passing.

% -----------------------------------------------------------------
\subsection{Where the work comes from}
\label{sec:pf:tasks}

Orders arrive continuously. At each timestep, orders are drawn from a
distribution $P_{\mathrm{ord}}$, an order being a finite multiset of
SKU--quantity pairs, and each order is converted into one or more transport tasks
by a task-generation rule. A task is a triple
\begin{equation}
  \tau_j = (r_j,\, s_j,\, g_j), \qquad
  r_j \in \mathbb{N}_0,\quad s_j \in V_{\mathrm{str}},\quad g_j \in V_{\mathrm{del}},
  \label{eq:task}
\end{equation}
consisting of a release time, a pickup vertex and a delivery vertex. The expected
number of tasks released per timestep is $\lambda_{\mathrm{task}}$.

The pickup vertex is not free. For the requested SKU $k_j$ it must satisfy
\begin{equation}
  x_{r_j}(k_j,\, s_j) > 0 .
  \label{eq:coupling}
\end{equation}
Equation~\eqref{eq:coupling} is the formal link between the two layers and the
central object of this thesis: it makes the spatial distribution of pickup
vertices a function of the storage configuration, so any change to $x_t$ changes
the routing problem that the controller subsequently faces. If inventory
depletion and replenishment are not simulated, $x_t$ is interpreted as a
placement template rather than a physical inventory level, and this
interpretation is stated with the results.

A task passes through three disjoint sets over its life. It enters the waiting
set $\mathcal{Q}_t$ at release, moves to the active set $\mathcal{B}_t$ when
execution begins, and enters the completed set $\mathcal{C}_t$ on delivery. One
convention for the boundary between $\mathcal{Q}_t$ and $\mathcal{B}_t$ is used
throughout: execution begins at \emph{assignment}, not at arrival of the agent at
$s_j$. Writing $d_j$ for the completion time, the service time is
\begin{equation}
  \zeta_j = d_j - r_j ,
  \label{eq:servicetime}
\end{equation}
which spans both phases and therefore cannot be improved by delaying assignment.
Figure~\ref{fig:lifecycle} shows the lifecycle and the interval that
\eqref{eq:servicetime} measures.

\begin{figure}[t]
  \centering
  \includegraphics[width=\linewidth]{Images/lifecycle.png}
  \caption{Lifecycle of a task. The service time $\zeta_j$ spans both the waiting
  and the active phase, so a controller cannot improve it by delaying assignment.
  The membership sets $\mathcal{Q}_t$, $\mathcal{B}_t$ and $\mathcal{C}_t$ are
  disjoint at every timestep, a property asserted at runtime in
  Section~\ref{sec:impl:validation}.}
  \label{fig:lifecycle}
\end{figure}

% -----------------------------------------------------------------
\subsection{Who decides}
\label{sec:pf:controllers}

A controller $\pi \in \Pi$ performs two logically distinct jobs. \emph{Task
assignment} maps free agents and waiting tasks to agent--task pairs; \emph{routing}
maps the resulting agent targets and the available system information to actions
$u_i(t) \in \mathcal{U}(\ell_i(t))$. Both jobs are treated as part of $\pi$, so
that a path planner is not described as though it also solves assignment.

Controllers differ in \emph{what information they may condition on}, and that is
the only distinction the problem formulation needs. Three information regimes are
studied and are summarised in Table~\ref{tab:information}. A centralised
controller may condition on the global state. A section-based controller
partitions the movable vertices into zones,
\begin{equation}
  \mathcal{Y} = \{Y_1,\dots,Y_p\}, \qquad
  \bigcup_{q=1}^{p} Y_q = V_{\mathrm{mov}}, \qquad
  Y_q \cap Y_{q'} = \emptyset \;\; (q \neq q'),
  \label{eq:partition}
\end{equation}
and may condition only on zone-local information, plus an explicit protocol at
zone boundaries. A decentralised controller conditions only on the local
observation defined in Section~\ref{sec:pf:observations} and selects actions from
a shared parameterised policy,
\[
  u_i(t) \sim \pi_\theta\bigl(\cdot \mid o_i(t)\bigr).
\]
Parameter sharing is appropriate for homogeneous agents and does not make
execution centralised; where global information is used during training only, the
method is described as centralised training with decentralised execution.

\begin{table}[t]
  \centering
  \caption{Information regimes. The three architectures are distinguished at the
  problem level only by what each component may condition on. Their algorithmic
  realisations are given in Section~\ref{sec:method:controllers}.}
  \label{tab:information}
  \begin{tabular}{lll}
    \toprule
    Architecture & Assignment conditions on & Routing conditions on \\
    \midrule
    Centralised     & global state & global state \\
    Section-based   & zone-local state, boundary protocol & zone-local state \\
    Decentralised   & global state (held fixed across arms) & $o_i(t)$ only \\
    \bottomrule
  \end{tabular}
\end{table}

Task assignment is held identical across the three architectures. Only routing
differs. This is a design decision of the study rather than a property of the
problem, and it is made so that a difference in results is attributable to the
routing architecture and not confounded with assignment quality.

% -----------------------------------------------------------------
\subsection{What a decentralised agent sees}
\label{sec:pf:observations}

A decentralised agent cannot see the warehouse. It sees a bounded neighbourhood
of its own position, and it can exchange information with nearby agents. Both
bounds are properties of the problem, because they define what a solution is
permitted to depend on.

For observation depth $d$, the visible vertex set and the induced subgraph are
\begin{equation}
  V^{(d)}_i(t) = \{ v \in V_{\mathrm{mov}} : d_G(\ell_i(t), v) \le d \},
  \qquad
  \mathcal{G}^{(d)}_i(t) = G\bigl[V^{(d)}_i(t)\bigr].
  \label{eq:localsubgraph}
\end{equation}
The size of \eqref{eq:localsubgraph} is bounded independently of $|V|$, which is
what allows a single policy to be evaluated on warehouses larger than those it
was trained on. The price is that anything beyond $d$ hops is invisible, so an
agent may commit to a move that becomes a conflict a few steps later.

Direction to a distant target is restored inside the local window by annotating
each visible vertex with its distance to the agent's current target $q_i(t)$,
\begin{equation}
  \eta_i(v,t) = d_G\bigl(v,\, q_i(t)\bigr).
  \label{eq:potential}
\end{equation}
The annotation \eqref{eq:potential} is a potential: its difference along an edge,
$\eta_i(w,t) - \eta_i(v,t)$, is negative exactly on shortest-path moves.
Following the negative gradient is the shortest path, and congestion avoidance is
precisely the choice to accept a non-negative step now in exchange for a lower
expected total later. Although $d_G$ may be precomputed, $\eta_i$ itself changes
whenever the target changes, which happens at every reassignment.

Agents within communication range form a second, time-varying graph whose nodes
are agents rather than locations,
\begin{equation}
  G^{\mathcal{A}}_t = (\mathcal{A}, E^{\mathcal{A}}_t), \qquad
  (a_i, a_j) \in E^{\mathcal{A}}_t \iff
  i \neq j \;\wedge\; d_G(\ell_i(t), \ell_j(t)) \le r_{\mathrm{com}} .
  \label{eq:commgraph}
\end{equation}
The environment graph $G$ and the communication graph $G^{\mathcal{A}}_t$ have
different node sets, different cardinalities and different edge semantics
(``can move to'' versus ``can exchange information with''); they must not be
conflated. The full observation available to agent $a_i$ is
\begin{equation}
  o_i(t) = \bigl( \mathcal{G}^{(d)}_i(t),\; \eta_i(\cdot, t),\;
                  \text{messages from } \{a_j : (a_i,a_j) \in E^{\mathcal{A}}_t\} \bigr).
  \label{eq:observation}
\end{equation}
How $o_i(t)$ is encoded into a vector, and what a message contains, are modelling
choices specified in Section~\ref{sec:method:model}.
Figure~\ref{fig:observation} illustrates both bounds on one instance.

\begin{figure}[t]
  \centering
  \includegraphics[width=\linewidth]{Images/observation.png}
  \caption{The observation of a single agent. The depth-$d$ induced subgraph is
  drawn at full strength and the rest of the warehouse is dimmed; numbers are the
  target distances $\eta_i(v,t)$. The dashed circle is the communication radius
  $r_{\mathrm{com}}$. One agent lies outside both bounds and is therefore
  invisible, although it may cause a conflict within a few timesteps.}
  \label{fig:observation}
\end{figure}

% -----------------------------------------------------------------
\subsection{The learning problem}
\label{sec:pf:game}

Because each agent acts from $o_i(t)$ while the environment evolves under the
joint action of all agents, the decentralised setting is a partially observable
stochastic game
\begin{equation}
  \mathcal{M} = \bigl( \mathcal{A},\, \mathcal{S},\, \{\mathcal{U}_i\}_{i=1}^{m},\,
                       P,\, \{\mathcal{O}_i\}_{i=1}^{m},\, \{O_i\}_{i=1}^{m},\,
                       \{R_i\}_{i=1}^{m},\, \gamma \bigr),
  \label{eq:posg}
\end{equation}
where $\mathcal{S}$ is the state space, $P$ the transition kernel,
$\mathcal{O}_i$ the observation space, $O_i$ the observation function, and
$R_i : \mathcal{S} \times \mathcal{U}^m \times \mathcal{S} \to \mathbb{R}$ the
reward function. The distinction between spaces and functions on the one hand and
realised states, observations and rewards on the other is maintained throughout.

A local congestion feature, defined so that an agent does not count itself, is
\begin{equation}
  \delta_i(t) =
  \frac{\bigl|\{ a_j \in \mathcal{A}\setminus\{a_i\} :
                 \ell_j(t) \in V^{(r_{\mathrm{cng}})}_i(t) \}\bigr|}
       {\max\bigl\{1,\; |V^{(r_{\mathrm{cng}})}_i(t)| - 1\bigr\}} .
  \label{eq:congestion}
\end{equation}
Because the agents are homogeneous and cooperative and share one parameter
vector, the learning objective is stated at team level,
\begin{equation}
  \max_{\theta} J(\theta)
    = \max_{\theta} \mathbb{E}_{\pi_\theta}
      \Bigl[ \sum_{t=0}^{T-1} \gamma^{t}\, \frac{1}{m} \sum_{i=1}^{m} R_i(t) \Bigr].
  \label{eq:objective}
\end{equation}
The reward functions $R_i$ are left unspecified here. A reward is an instrument
for obtaining a policy that scores well on the system-level criterion of
Section~\ref{sec:pf:measures}; it can be redesigned without changing the problem,
and it is therefore specified in Section~\ref{sec:method:rl}.

% -----------------------------------------------------------------
\subsection{How the two layers close the loop}
\label{sec:pf:coupling}

Let $s_t \in \mathcal{S}$ contain the agent states, the task lifecycle sets and
the current storage configuration. The system runs on two clocks. On every
timestep, the fast dynamics advance under the controller,
\[
  s_{t+1} \sim P(\cdot \mid s_t, u_t), \qquad u_t \sim \pi(\cdot \mid s_t),
\]
with the corresponding product of local policies under decentralised execution.
At storage epochs only, the slow dynamics additionally apply
\eqref{eq:storageupdate}.

The loop closes in two directions. Forward, a change in $x_t$ changes which
vertices appear as pickups through the coupling condition
\eqref{eq:coupling}, and therefore changes the routing problem. Backward, the
realised routing statistics $\mu_t$ and $w_t$ enter the next storage update
through \eqref{eq:storageupdate}. Neither direction exists in a study that fixes
one layer, and Figure~\ref{fig:loop} shows that the fixed-storage baseline is
exactly this diagram with the backward path removed.

\begin{figure}[t]
  \centering
  \includegraphics[width=\linewidth]{Images/loop.png}
  \caption{The two-timescale loop. The fast path runs every timestep; the slow
  feedback path runs every $\Delta$ steps. Deleting the feedback path yields the
  fixed-storage baseline, so the experimental contrast of this thesis is a
  contrast between this diagram with and without one arrow.}
  \label{fig:loop}
\end{figure}

The separation of timescales has a consequence that constrains the study. Between
epochs, the fast process is approximately stationary and the storage rule sees
only its aggregate statistics. If $\Delta$ is too small relative to the mixing
time of the fast process, $F$ reacts to noise and the task distribution becomes
non-stationary from the learner's point of view; if $\Delta$ is too large, the
estimates $\hat{\rho}, \hat{\mu}, \hat{w}$ are stale. The choice of $\Delta$ is
therefore an experimental variable, not a constant.

% -----------------------------------------------------------------
\subsection{How a run is scored}
\label{sec:pf:measures}

A run is scored on completed work, on the cost of producing it, and on how the
traffic was distributed while producing it. For the completed set $\mathcal{C}_T$
over a horizon $T$, throughput and mean service time are
\begin{equation}
  \Lambda_T = \frac{|\mathcal{C}_T|}{T}, \qquad
  \bar{\zeta}_T = \frac{1}{|\mathcal{C}_T|}
                  \sum_{\tau_j \in \mathcal{C}_T} (d_j - r_j),
  \label{eq:throughput}
\end{equation}
and the movement cost per completed task is
\begin{equation}
  \bar{c}_T = \frac{\sum_{i=1}^{m} \sum_{t=0}^{T-1}
                    \hat{c}\bigl(\ell_i(t), \ell_i(t+1)\bigr)}
                   {|\mathcal{C}_T|} .
  \label{eq:movementcost}
\end{equation}
Runs with $|\mathcal{C}_T| = 0$ are reported separately rather than assigned an
arbitrary ratio. Because a completion-only statistic can improve while the system
falls behind, the backlog $|\mathcal{Q}_T| + |\mathcal{B}_T|$ and the age
distribution of unfinished tasks are always reported alongside
\eqref{eq:throughput}.

Traffic distribution is measured separately from delay. Let $\mu_T(e)$ be the
number of traversals of edge $e$ over the evaluation window,
$E^{+}_T = \{e \in E : \mu_T(e) > 0\}$, and
$p_T(e) = \mu_T(e) / \sum_{e' \in E} \mu_T(e')$. For $|E^{+}_T| > 1$ the
normalised entropy is
\begin{equation}
  H_T = -\frac{1}{\log |E^{+}_T|} \sum_{e \in E^{+}_T} p_T(e) \log p_T(e)
        \;\in [0,1],
  \label{eq:entropy}
\end{equation}
so $H_T = 1$ means perfectly uniform edge usage and $H_T \to 0$ means all traffic
on one edge. Since entropy measures dispersion, the concentration index is its
complement,
\begin{equation}
  C_T = 1 - H_T .
  \label{eq:concentration}
\end{equation}
The direction matters: adding $+H_T$ to a minimised cost would reward funnelling
all agents down one aisle. Two caveats accompany \eqref{eq:entropy}. The
normalisation depends on $|E^{+}_T|$, so runs using very different numbers of
edges are not strictly comparable; and concentration is not congestion, because a
concentrated flow at low load causes no delay. Congestion and concentration are
therefore reported as separate outcomes.
Figure~\ref{fig:concentration} illustrates the two extremes.

\begin{figure}[t]
  \centering
  \includegraphics[width=\linewidth]{Images/concentration.png}
  \caption{Two edge-usage patterns on the same graph, drawn on a shared width
  scale, with the resulting $H_T$ and $C_T$. The panels are schematic
  constructions that define the measure; measured values from experimental runs
  are reported in Section~\ref{sec:experiments}.}
  \label{fig:concentration}
\end{figure}

The system-level criterion aggregates the minimisation-oriented measures,
\begin{equation}
  \mathcal{J}(F,\pi) = q_{\mathrm{svc}}\, \mathbb{E}[\bar{\zeta}_T]
                     + q_{\mathrm{mov}}\, \mathbb{E}[\bar{c}_T]
                     + q_{\mathrm{wait}}\, \mathbb{E}[W_T]
                     + q_{\mathrm{conc}}\, \mathbb{E}[C_T]
                     + q_{\mathrm{back}}\, \mathbb{E}[B_T],
  \label{eq:functional}
\end{equation}
with all $q_\bullet \ge 0$, $B_T = |\mathcal{Q}_T| + |\mathcal{B}_T|$, and $W_T$
a waiting measure whose operational definition is given in
Section~\ref{sec:impl:logging}. Throughput is deliberately excluded from
\eqref{eq:functional}: it is maximisation-oriented and, under a fixed arrival
rate, largely redundant with $\bar{\zeta}_T$, so including it with a negative
weight would double-count. It is reported separately.

The formal decision problem is
\begin{equation}
  (F^{\star}, \pi^{\star}) \in
  \arg\min_{F \in \mathcal{F},\, \pi \in \Pi} \mathcal{J}(F,\pi)
  \quad \text{subject to \eqref{eq:feasiblestorage}, \eqref{eq:transition},
  \eqref{eq:vertexconflict}, \eqref{eq:swapconflict} and \eqref{eq:coupling}.}
  \label{eq:decisionproblem}
\end{equation}
Equation~\eqref{eq:decisionproblem} is a framing device. This thesis does not
solve it globally; it estimates the effect of selected $(F,\pi)$ combinations
under a factorial design. Concretely, the empirical question is whether there
exists an adaptive rule $F$ such that
\begin{equation}
  \mathcal{J}(F, \pi) \;<\; \mathcal{J}(F_{\mathrm{fix}}, \pi)
  \quad \text{by a margin exceeding seed variance,}
  \label{eq:rqformal}
\end{equation}
and whether the size of that margin depends on the information regime of $\pi$
given in Table~\ref{tab:information} and on the load parameters $m$ and
$\lambda_{\mathrm{task}}$. Section~\ref{sec:method} realises every symbol
appearing in \eqref{eq:rqformal}.

% -----------------------------------------------------------------
\subsection{Scope and assumptions}
\label{sec:pf:scope}

The model above is deliberately narrower than a physical warehouse, and the
following are not represented. Agents are homogeneous, occupy exactly one vertex,
and carry at most one task; there is no capacity, battery, or kinematic model, so
turning and acceleration costs are absorbed into $c(e)$. Execution is
deterministic: an accepted action always succeeds, so the transition kernel $P$
is stochastic only through order arrivals. Human pickers and manual interventions
are absent. Inventory relocation is abstract --- \eqref{eq:storageupdate} changes
the placement used to generate future tasks, but the physical movement of stock
is not simulated, and if $\chi = 0$ in \eqref{eq:storageobjective} that movement
is also free. Replenishment and depletion are modelled only if stated, and
otherwise $x_t$ is a placement template.

Each of these exclusions bounds the claims that the experiments can support. In
particular, results about the value of congestion feedback are conditional on
relocation being cheap, and results about decentralised control are conditional
on deterministic execution.
\end{document}